\Needspace{12\baselineskip}
\originalchapter{计数与能量分辨}
\originalsection{Poisson计数和死时间}
对稳定、独立且小时间内双事件概率可忽略的过程，时间 $T$ 内计数均值 $\mu=RT$。将 $T$ 分成大量小区间，各区间事件概率 $R\Delta t$，二项分布极限给 $P(n)=\e^{-\mu}\mu^n/n!$，均值和方差都等于 $\mu$。Junior Lab公开模块使用闪烁探测器及放大、阈值和计数链，比较不同计数率下等时段分布，并与随机模拟比较。

阈值过低会收入电子学噪声，过高会漏掉真实脉冲；改变高压或增益也改变脉冲幅值分布。计数率低时概率明显不对称，不能只凭一条对称高斯曲线判断异常。高计数时高斯近似才逐渐适用。

非延长型死时间模型中，每记录一次事件后有固定时间 $\tau_d$ 不能记录，之后再等待平均 $R^{-1}$。平均记录率 $r=1/(\tau_d+R^{-1})=R/(1+R\tau_d)$，反推 $R=r/(1-r\tau_d)$。这一模型与每次未记录事件都延长失效时间的延长型模型不同。死时间还会改变计数方差，因此“均值方差不等”不一定说明放射源有物理关联。

有背景时分别测量源加背景 $n_s$、背景 $n_b$，时间 $T_s,T_b$，净率估计 $\hat R=n_s/T_s-n_b/T_b$，独立Poisson近似给方差 $n_s/T_s^2+n_b/T_b^2$。负净率可由涨落产生，不能把原始负估计删去后再进行无偏统计。

\originalsection{Compton散射的几何和校准}
光子初能量 $E$，与近似静止电子散射后能量 $E'$，偏转角 $\theta$。能动量守恒将电子动量消去：$m_ec^2(E-E')=EE'(1-\cos\theta)$，因而
\[
\frac1{E'}-\frac1E=\frac{1-\cos\theta}{m_ec^2},\qquad
E'=\frac E{1+(E/m_ec^2)(1-\cos\theta)}.
\]
Junior Lab手册选661.6 keV伽马光子及能谱探测。可从散射光谱峰随角度变化检验关系，也可比较反冲电子能量 $E-E'$。实际源、散射体和探测器大小使接受角有分布，有限角宽不能简单以中心角代入而忽略展宽。

多道分析器首先要用已知峰标定通道和能量，检查线性、峰宽和背景。谱中的Compton连续部分、多次散射和全能峰不是同一对象。峰位置应通过包含背景的模型拟合，不能把最大计数通道当作无不确定度的峰心。能量标定参数的协方差应传到每个散射角的结果；各峰共享同一标定，残差因此可能相关。

\originalsection{Franck--Hertz碰撞谱}
加热阴极发射电子，经栅极电压加速后，以微小反向电势筛选仍有足够动能抵达阳极的电子。汞蒸气中，电子能量低于激发阈值时主要发生近弹性碰撞；超过第一激发能时能将能量交给原子，自己减速，阳极电流出现谷。继续增大电压允许多次激发，谷间距近似给 $\Delta E/e$，汞的对应尺度约4.9 V。

不能用第一个谷的绝对电压直接求激发能，因为阴极和栅极接触电势、反向电势及能量分布会移动整体位置。拟合谷序号 $j$ 与电压 $V_j=V_0+j\Delta V$，斜率比首谷位置更稳健。温度调节蒸气密度，改变平均自由程，过低温使非弹性碰撞不充分，过高温使碰撞过多、结构模糊。

手册还讨论空间电荷限制发射、汞电离以及Ramsauer--Townsend弹性散射现象；这些不是同一阈值机制。Ramsauer散射截面随能量的极小值依赖电子波与原子势的散射，不可解释成“原子不允许吸收这份能量”。本卷解释其区别，未把不可获的课前网页问题恢复为原题。

\originalsection{Poisson似然和背景估计}
对等时段计数 $n_1,\ldots,n_M$，对数似然 $\ell(\mu)=\sum[n_i\ln\mu-\mu-\ln(n_i!)]$，求导给 $\hat\mu=\bar n$，Fisher信息 $M/\mu$，大样本标准误 $\sqrt{\mu/M}$。这由概率模型推出，而不是任意设置误差为平均值的平方根。对很低总计数，置信区间应由离散概率反演或明确的似然区间求出；观察零计数也有有限上界，不意味着真率精确为零。

源加背景与独立背景两次测量的净计数差服从两Poisson差的分布，净率估计可以为负。若要约束物理源率非负，应在原始两个Poisson模型中引入共同背景率求似然，而不是先截断净数据。改变测背景时间会改变误差，在固定总时间下最优分配依赖源加背景和背景的计数率。

计数方差大于平均可能来自源强漂移：若条件均值为随机 $\mu$，全方差公式给 $\operatorname{Var}n=\avg\mu+\operatorname{Var}\mu$。因此过散布既可能是事件关联，也可能是慢变的仪器或源。查看时间序列并按短窗口重复分析有助于区分。

\originalsection{死时间模型的推导和鉴别}
非延长型过程每条记录的间隔是固定失活 $\tau_d$ 加一个指数等待，因此记录间隔均值 $\tau_d+1/R$、方差 $1/R^2$。长期记录数的Fano因子约 $\operatorname{Var}n/\avg n=1/(1+R\tau_d)^2$；它本身可小于1。延长型过程每次到达都会重启失活，事件能被记录要求之前 $\tau_d$ 内无到达，故 $r=R\e^{-R\tau_d}$。后者在高真率时记录率会下降，多值反推问题也随之出现。

因此应以不同源强、双源法或脉冲时间间隔检验死时间。两个源的合并记录率一般不是各自记录率的和。附带事件级模拟采用非延长型，$R=200\,\mathrm{s^{-1}},\tau_d=.002$ s，2000 s记录率 $142.8875\,\mathrm{s^{-1}}$，理论 $142.8571\,\mathrm{s^{-1}}$，只证明所写模型代码与更新过程相符。

\originalsection{Compton关系的能动量推导}
初光子动量 $E/c$，末光子 $E'/c$，角度为 $\theta$。电子末动量平方为 $p_e^2c^2=E^2+E'^2-2EE'\cos\theta$。电子总能 $m_ec^2+E-E'$，代入电子质量壳关系并消去 $E^2,E'^2$，得到 $2m_ec^2(E-E')=2EE'(1-\cos\theta)$。用波长 $E=hc/\lambda$ 可写 $\lambda'-\lambda=(h/(m_ec))(1-\cos\theta)$。直接从守恒律出发解释了静止自由电子假设的位置。

散射角的不确定度传到能量为 $u_{E',\theta}=|\partial_\theta E'|u_\theta$，$\partial_\theta E'=-E\alpha\sin\theta/[1+\alpha(1-\cos\theta)]^2$、$\alpha=E/(m_ec^2)$。有限接受角造成的是分布卷积，而非只有一个误差条；在大曲率区，平均峰位也不等于中心角预测。束缚电子的初始动量产生额外展宽，多个散射体事件也可改变谱。

若全能峰用Gaussian加线性背景 $\mu(C)=A\exp[-(C-C_0)^2/(2s^2)]+b_0+b_1C$ 拟合，应按各通道Poisson似然求参数。换拟合窗口、背景阶数并检查峰肩，可估计模型依赖。峰宽和峰心标准误不同：很多计数可使峰心很准，即使单事件分辨率并不窄。

\originalsection{Franck--Hertz温度、平均自由程与峰间距}
电子碰撞概率取决于蒸气密度 $n$、能量相关截面 $\sigma(E)$ 和传播距离，平均自由程 $\ell(E)=1/[n\sigma(E)]$。调炉温改变 $n$，从而决定电子在跨栅极前能否多次得到阈值能量并碰撞。若电势在长度 $d_g$ 上近似线性，得到激发能所需距离约 $x_{\rm exc}=d_g\Delta E/(eV)$，这帮助解释为何谷序列随温度及电压出现，而不是仅有一个阈值。

谷位置可受多次碰撞、电子初能分布、空间电荷和接触势共同影响。以多个谷序号直线拟合求 $\Delta V$，并比较不同温度和不同识别窗口，比只挑相邻一对谷更稳健。最初几个谷不等距时要展示原曲线，说明是否进入预期渐近，而不能悄悄删到只剩理想等距点。

实验显示原子能交换具有特定激发尺度，但峰间距并非所有原子能级的完整谱。汞第一激发、后续电离和弹性散射结构对应不同物理过程，课堂报告应分别给阈值、装置条件和识别依据。
